% Audit-summary figure of the ERL Letter (D-100 display item 3). For now the Phase A audit
% figure (figures/fig4_audit.py), by the lead's instruction of 2026-10-21; the lead decides
% whether to slim it down. Its type is still the Phase A size (5-7 pt), below ERL's 8 pt
% minimum, so it must be redrawn before submission (D-107). Legend adapted from
% figure_captions_draft.tex (fig:audit) to ERL style; every number is a \metric{} macro.
\begin{figure}[!htbp]
  \centering
  \includegraphics[width=\linewidth]{../figures/out/fig4_audit.pdf}
  \caption{\textbf{Auditing the bathymetry grids behind emergence projections.}
  (a) Area bias of the sea that each global bathymetry grid\citep{gebco2026,gebco2022,etopo2022,macferrin2025,tozer2019,amante2009} predicts at an observed level, against the observed sea, for each observed level of each epoch inside the study domain ($n = \metric{h1.n_units_domain}$ levels with observed sea: \metric{h1.n_units_domain.E1} on the rising limb before 1995, \metric{h1.n_units_domain.E2} in 1995--2010 and \metric{h1.n_units_domain.E3} after 2010), as a percentage of the observed sea area in the cells observed at that level; positive, more sea predicted than observed. Upper, each grid shifted by the single vertical constant that best fits it (GEBCO\_2026 \metric{h1.shape_offset.gebco2026_m}~m, GEBCO\_2022 \metric{h1.shape_offset.gebco2022_m}~m, ETOPO 2022 \metric{h1.shape_offset.etopo2022_m}~m, SRTM15+ \metric{h1.shape_offset.srtm15p_m}~m, ETOPO1 \metric{h1.shape_offset.etopo1_m}~m); lower, each grid on the documented datum chain without a constant. Line style, epoch.
  (b) Shoreline displacement by level band, after the constant: per band (edges \metric{h1.disp_band.edge_0_m}, \metric{h1.disp_band.edge_1_m}, \metric{h1.disp_band.edge_2_m}, \metric{h1.disp_band.edge_3_m}, \metric{h1.disp_band.edge_4_m} and \metric{h1.disp_band.edge_5_m}~m on the altimetry reference of the Database for Hydrological Time Series of Inland Waters), the median over sectors and levels of the per-level median distance between predicted and observed shorelines. Dashed line, \metric{h1.disp_band.edge_2_m}~m.
  (c) Median absolute shoreline displacement per \metric{params.block_size_km}~km block, pooled over all observed levels, after the constant, for GEBCO\_2026 and ETOPO 2022; the other grids and the grids without the constant are in the Supplementary Information (section~\SIaudit). Blocks with fewer than \metric{params.h1_min_block_n} distances are not drawn; the number drawn is printed on each map. Light blue, open sea in \metric{params.lowstand_label}; Lambert azimuthal equal-area projection centred at \metric{params.laea_lat0_deg}\textdegree{}N, \metric{params.laea_lon0_deg}\textdegree{}E.
  (d) Per sector, the median over the evaluable levels of each grid's per-level median displacement, after the constant, against the sector's observation noise floor (grey bar: the median over the same levels of the 95th percentile of the distances between the shorelines of two independent halves of the observations). Filled, the grid mispredicts the sector (more than half of its evaluable levels exceed the floor; pre-registered rule); open, it does not.}
  \label{fig:erl_audit}
\end{figure}
